\documentclass[a4paper,12pt,twoside]{memoir}

% Castellano
\usepackage[spanish,es-tabla]{babel}
\selectlanguage{spanish}
\usepackage[utf8]{inputenc}
\usepackage[T1]{fontenc}
\usepackage{lmodern} % Scalable font
\usepackage{microtype}
\usepackage{placeins}
% Bibliography management
\usepackage[backend=bibtex]{biblatex}
\addbibresource{bibliografia.bib}

\RequirePackage{booktabs}
\RequirePackage[table]{xcolor}
\RequirePackage{xtab}
\RequirePackage{multirow}

% Links
\usepackage[colorlinks]{hyperref}
\hypersetup{
	allcolors = {red}
}



% Ecuaciones
\usepackage{amsmath}

% Rutas de fichero / paquete
\newcommand{\ruta}[1]{{\sffamily #1}}

% Párrafos
\nonzeroparskip


% Imagenes
\usepackage{graphicx}
\newcommand{\imagen}[2]{
	\begin{figure}[!h]
		\centering
		\includegraphics[width=0.9\textwidth]{#1}
		\caption{#2}\label{fig:#1}
	\end{figure}
	\FloatBarrier
}

\newcommand{\imagenflotante}[2]{
	\begin{figure}%[!h]
		\centering
		\includegraphics[width=0.9\textwidth]{#1}
		\caption{#2}\label{fig:#1}
	\end{figure}
}



% El comando \figura nos permite insertar figuras comodamente, y utilizando
% siempre el mismo formato. Los parametros son:
% 1 -> Porcentaje del ancho de página que ocupará la figura (de 0 a 1)
% 2 --> Fichero de la imagen
% 3 --> Texto a pie de imagen
% 4 --> Etiqueta (label) para referencias
% 5 --> Opciones que queramos pasarle al \includegraphics
% 6 --> Opciones de posicionamiento a pasarle a \begin{figure}
\newcommand{\figuraConPosicion}[6]{%
  \setlength{\anchoFloat}{#1\textwidth}%
  \addtolength{\anchoFloat}{-4\fboxsep}%
  \setlength{\anchoFigura}{\anchoFloat}%
  \begin{figure}[#6]
    \begin{center}%
      \Ovalbox{%
        \begin{minipage}{\anchoFloat}%
          \begin{center}%
            \includegraphics[width=\anchoFigura,#5]{#2}%
            \caption{#3}%
            \label{#4}%
          \end{center}%
        \end{minipage}
      }%
    \end{center}%
  \end{figure}%
}

%
% Comando para incluir imágenes en formato apaisado (sin marco).
\newcommand{\figuraApaisadaSinMarco}[5]{%
  \begin{figure}%
    \begin{center}%
    \includegraphics[angle=90,height=#1\textheight,#5]{#2}%
    \caption{#3}%
    \label{#4}%
    \end{center}%
  \end{figure}%
}
% Para las tablas
\newcommand{\otoprule}{\midrule [\heavyrulewidth]}
%
% Nuevo comando para tablas pequeñas (menos de una página).
\newcommand{\tablaSmall}[5]{%
 \begin{table}
  \begin{center}
   \rowcolors {2}{gray!35}{}
   \begin{tabular}{#2}
    \toprule
    #4
    \otoprule
    #5
    \bottomrule
   \end{tabular}
   \caption{#1}
   \label{tabla:#3}
  \end{center}
 \end{table}
}

%
% Nuevo comando para tablas pequeñas (menos de una página).
\newcommand{\tablaSmallSinColores}[5]{%
 \begin{table}[H]
  \begin{center}
   \begin{tabular}{#2}
    \toprule
    #4
    \otoprule
    #5
    \bottomrule
   \end{tabular}
   \caption{#1}
   \label{tabla:#3}
  \end{center}
 \end{table}
}

\newcommand{\tablaApaisadaSmall}[5]{%
\begin{landscape}
  \begin{table}
   \begin{center}
    \rowcolors {2}{gray!35}{}
    \begin{tabular}{#2}
     \toprule
     #4
     \otoprule
     #5
     \bottomrule
    \end{tabular}
    \caption{#1}
    \label{tabla:#3}
   \end{center}
  \end{table}
\end{landscape}
}

%
% Nuevo comando para tablas grandes con cabecera y filas alternas coloreadas en gris.
\newcommand{\tabla}[6]{%
  \begin{center}
    \tablefirsthead{
      \toprule
      #5
      \otoprule
    }
    \tablehead{
      \multicolumn{#3}{l}{\small\sl continúa desde la página anterior}\\
      \toprule
      #5
      \otoprule
    }
    \tabletail{
      \hline
      \multicolumn{#3}{r}{\small\sl continúa en la página siguiente}\\
    }
    \tablelasttail{
      \hline
    }
    \bottomcaption{#1}
    \rowcolors {2}{gray!35}{}
    \begin{xtabular}{#2}
      #6
      \bottomrule
    \end{xtabular}
    \label{tabla:#4}
  \end{center}
}

%
% Nuevo comando para tablas grandes con cabecera.
\newcommand{\tablaSinColores}[6]{%
  \begin{center}
    \tablefirsthead{
      \toprule
      #5
      \otoprule
    }
    \tablehead{
      \multicolumn{#3}{l}{\small\sl continúa desde la página anterior}\\
      \toprule
      #5
      \otoprule
    }
    \tabletail{
      \hline
      \multicolumn{#3}{r}{\small\sl continúa en la página siguiente}\\
    }
    \tablelasttail{
      \hline
    }
    \bottomcaption{#1}
    \begin{xtabular}{#2}
      #6
      \bottomrule
    \end{xtabular}
    \label{tabla:#4}
  \end{center}
}

%
% Nuevo comando para tablas grandes sin cabecera.
\newcommand{\tablaSinCabecera}[5]{%
  \begin{center}
    \tablefirsthead{
      \toprule
    }
    \tablehead{
      \multicolumn{#3}{l}{\small\sl continúa desde la página anterior}\\
      \hline
    }
    \tabletail{
      \hline
      \multicolumn{#3}{r}{\small\sl continúa en la página siguiente}\\
    }
    \tablelasttail{
      \hline
    }
    \bottomcaption{#1}
  \begin{xtabular}{#2}
    #5
   \bottomrule
  \end{xtabular}
  \label{tabla:#4}
  \end{center}
}



\definecolor{cgoLight}{HTML}{EEEEEE}
\definecolor{cgoExtralight}{HTML}{FFFFFF}

%
% Nuevo comando para tablas grandes sin cabecera.
\newcommand{\tablaSinCabeceraConBandas}[5]{%
  \begin{center}
    \tablefirsthead{
      \toprule
    }
    \tablehead{
      \multicolumn{#3}{l}{\small\sl continúa desde la página anterior}\\
      \hline
    }
    \tabletail{
      \hline
      \multicolumn{#3}{r}{\small\sl continúa en la página siguiente}\\
    }
    \tablelasttail{
      \hline
    }
    \bottomcaption{#1}
    \rowcolors[]{1}{cgoExtralight}{cgoLight}

  \begin{xtabular}{#2}
    #5
   \bottomrule
  \end{xtabular}
  \label{tabla:#4}
  \end{center}
}


















\graphicspath{ {./img/} }

% Capítulos
\chapterstyle{bianchi}
\newcommand{\capitulo}[2]{
	\setcounter{chapter}{#1}
	\setcounter{section}{0}
	\chapter*{#2}
	\addcontentsline{toc}{chapter}{#2}
	\markboth{#2}{#2}
}

% Apéndices
\renewcommand{\appendixname}{Apéndice}
\renewcommand*\cftappendixname{\appendixname}

\newcommand{\apendice}[1]{
	%\renewcommand{\thechapter}{A}
	\chapter{#1}
}

\renewcommand*\cftappendixname{\appendixname\ }

% Formato de portada
\makeatletter
\usepackage{xcolor}
\newcommand{\tutor}[1]{\def\@tutor{#1}}
\newcommand{\course}[1]{\def\@course{#1}}
\definecolor{cpardoBox}{HTML}{E6E6FF}
\def\maketitle{
  \null
  \thispagestyle{empty}
  % Cabecera ----------------
\noindent\includegraphics[width=\textwidth]{cabecera}\vspace{1cm}%
  \vfill
  % Título proyecto y escudo informática ----------------
  \colorbox{cpardoBox}{%
    \begin{minipage}{.8\textwidth}
      \vspace{.5cm}\Large
      \begin{center}
      \textbf{TFG del Grado en Ingeniería Informática}\vspace{.6cm}\\
      \textbf{\LARGE\@title{}}
      \end{center}
      \vspace{.2cm}
    \end{minipage}

  }%
  \hfill\begin{minipage}{.20\textwidth}
    \includegraphics[width=\textwidth]{escudoInfor}
  \end{minipage}
  \vfill
  % Datos de alumno, curso y tutores ------------------
  \begin{center}%
  {%
    \noindent\LARGE
    Presentado por \@author{}\\ 
    en Universidad de Burgos --- \@date{}\\
    Tutor: \@tutor{}\\
  }%
  \end{center}%
  \null
  \cleardoublepage
  }
\makeatother

\newcommand{\nombre}{Yi Peng Ji} %%% cambio de comando

% Datos de portada
\title{{\Huge UBUMonitor}\\[0.5cm]Monitorización de alumnos en la plataforma Moodle.}
\author{\nombre}
\tutor{Dr.D. Raúl Marticorena Sánchez\\y Dr.D. Carlos Pardo Aguilar}
\date{\today}

\begin{document}

\maketitle


\newpage\null\thispagestyle{empty}\newpage


%%%%%%%%%%%%%%%%%%%%%%%%%%%%%%%%%%%%%%%%%%%%%%%%%%%%%%%%%%%%%%%%%%%%%%%%%%%%%%%%%%%%%%%%
\thispagestyle{empty}


\noindent\includegraphics[width=\textwidth]{cabecera}\vspace{1cm}

\noindent Dr.D. Raúl Marticorena Sánchez y Dr.D. Carlos Pardo Aguilar, profesores del departamento de Ingeniería Civil, área de Lenguajes y Sistemas Informáticos.

\noindent Expone:

\noindent Que el alumno D. \nombre, con NIE X2149097T, ha realizado el Trabajo final de Grado en Ingeniería Informática titulado ``UBUMonitor - Monitorización de alumnos en la plataforma Moodle''. 

\noindent Y que dicho trabajo ha sido realizado por el alumno bajo la dirección del que suscribe, en virtud de lo cual se autoriza su presentación y defensa.

\begin{center} %\large
En Burgos, {\large \today}
\end{center}

\vfill\vfill\vfill

% Author and supervisor
\begin{minipage}{0.45\textwidth}
\begin{flushleft} %\large
Vº. Bº. del Tutor:\\[2cm]
Dr. Raúl Marticorena Sánchez
\end{flushleft}
\end{minipage}
\hfill
\begin{minipage}{0.45\textwidth}
\begin{flushleft} %\large
Vº. Bº. del co-tutor:\\[2cm]
Dr. Carlos Pardo Aguilar
\end{flushleft}
\end{minipage}
\hfill

\vfill

% para casos con solo un tutor comentar lo anterior
% y descomentar lo siguiente
%Vº. Bº. del Tutor:\\[2cm]
%D. nombre tutor


\newpage\null\thispagestyle{empty}\newpage




\frontmatter

% Abstract en castellano
\renewcommand*\abstractname{Resumen}
\begin{abstract}

El auge de las tecnologías ha alcanzado el ámbito educacional con las aulas de aprendizaje \textit{online} y enseñanza virtual con Moodle como uno de los protagonistas más importantes.

Este hecho ha transformado también el método de seguimiento de la evolución del alumnado durante el curso.

En Moodle permite el seguimiento mediante la visualización de las calificaciones  y los registros de acceso de los alumnos. Sin embargo la forma de mostrar las calificaciones y registros de acceso al curso de los alumnos en Moodle es bastante laxa y poco práctica. Este inconveniente ha motivado la creación del presente proyecto como solución.

UBUMonitor es una aplicación de escritorio hecho en Java orientado a docentes que usen los servicios de Moodle. Permite a los profesores, previa identificación, acceder a datos estadísticos diferentes sobre las calificaciones y registros de los alumnos y mostrarlos en distintos tipos de gráficas, una forma mucho más práctica y visual. 


\end{abstract}

\renewcommand*\abstractname{Descriptores}
\begin{abstract}
Moodle, Educational data mining, Visualización, JavaFX, Servicios web, JSON, ChartJs.
\end{abstract}

\clearpage

% Abstract en inglés
\renewcommand*\abstractname{Abstract}
\begin{abstract}
The technology boom has reached the educational field with online learning classrooms and virtual teaching with Moodle as one of the most important protagonists.

This fact has also transformed the method of monitoring the evolution of students during the course.

In Moodle, it allows students to monitor their grades and access records. However, the way in which grades and access records are displayed in Moodle is quite lax and impractical. This inconvenience has led to the creation of this project as a solution.

UBUMonitor is a Java desktop application for teachers using Moodle services. It allows teachers, after identification, to access different statistical data on student grades and records and display them in different types of graphs, a much more practical and visual way. 
\end{abstract}

\renewcommand*\abstractname{Keywords}
\begin{abstract}
Moodle, Educational data mining, Visualization, JavaFX, Web services, JSON, ChartJs.
\end{abstract}

\clearpage

% Indices
\tableofcontents

\clearpage

\listoffigures

\clearpage

\listoftables
\clearpage

\mainmatter
\capitulo{1}{Introducción}

El presente proyecto es una continuación de dos trabajos anteriores, UBUGrades 2.0\cite{nogal_santamaria_continuacion_2017} de Félix Nogal Santamaría y UBULogs\cite{armengol_aplicacion_2017} de Oscar Fernández Armengol. 

En ambos trabajos, la idea principal de la aplicación es el seguimiento de las calificaciones por una parte y registros de eventos por otra. Todo esto representado mediante diversas gráficas, una forma sencilla y visual de mostrar los datos.

El desarrollo de la nueva aplicación se ha integrado ambos proyectos usando UBUGrades como base, sin embargo, UBULogs no funcionaba correctamente debido a los cambios en Moodle. Este hecho ha motivado empezar de cero en esta parte siguiendo una vía alternativa.

Se ha centrado también en la refactorización del código, guardar los datos en ficheros cifrados y nuevos gráficos para los registros de accesos.

En resumidas cuentas, está nueva versión del proyecto permite visualizar en diferentes gráficas las calificaciones de los alumnos por una parte y sobre los registros de eventos por otra parte. Todo esto con medidas de almacenaje de datos en local, y cifrados al contener información sensible. Esto posibilita acceder a los datos sin tener que realizar peticiones de recuperación de datos al servidor constantemente cada vez que se use la aplicación.
\capitulo{2}{Objetivos del proyecto}

En este apartado mostraremos los diferentes objetivos que han surgido para realizar el proyecto.

\section{Objetivos generales}

\begin{itemize}
	\item Modificar el código para conseguir una mejor mantenibiliad y facilidad de añadir nuevas funcionalidades. 
	\item Guardar los datos de Moodle en ficheros cifrados con la finalidad de acceder más rápidamente a los datos sin tener que actualizar del servidor continuamente.
	\item Implementar la parte de los registros (\textit{logs}), la descarga de los ficheros CSV desde Moodle hasta su clasificación según su tipo de evento.
	\item Gráficas nuevas para mostrar el número de registros por cada usuario y evento.
	
\end{itemize}

\section{Objectivos técnicos}
\begin{itemize}
	\item Aprendizaje sobre las REST API, incluyendo los métodos de petición HTTP\cite{noauthor_http_nodate}.
	\item Uso de Git como método de control de versiones con la plataforma GitHub.
	\item Continuar el desarrollo con la librería JavaFx para la interfaz gráfica.
	\item Mejorar la mantenibiliad del software aplicando patrones de diseño.
\end{itemize}

\section{Objetivos personales}
\begin{itemize}
	\item Aplicar lo máximo posible de lo aprendido de las asignaturas de la universidad.
	\item Mejorar en la programación de Java profundizando en la genericidad, programación funcional y patrones de diseño.
	\item Aportar a la comunidad de Moodle ayudando a la extracción de información de calificaciones y registros de los alumnos.
	\item Mantener la motivación durante el desarrollo del trabajo, así es más fácil disfrutar realizar el proyecto.
\end{itemize}

\capitulo{3}{Conceptos teóricos}

En este apartado comentaremos sobre conceptos que aparecen en Moodle.



\section{Conceptos de Moodle}

//TODO Introduccion a moodle

\subsection{Módulo del curso}
El módulo del curso \cite{noauthor_course_nodate} (en inglés \textit{course module}) también abreviado como `cm', representa las \textbf{actividades} y \textbf{recursos} del curso (ejemplos en la figura \ref{fig:course_modules}). Estos contienen información sobre qué curso y sección se muestra, además de la configuración de visibilidad y otros datos relevantes. 

\imagenflotante{course_modules}{Ejemplos de módulos del curso que actividades y recursos.}

Los identificadores únicos de los módulos del curso se definen como \textbf{`cmid'} y cada tipo de módulo del curso posee un código nemotécnico.


\subsubsection{Códigos nemotécnicos de los tipos de módulos del curso}

Existen dos formas de buscar el código nemotécnico de un módulo del curso:
\begin{itemize}
	\item A través de la función de servicio web \textbf{\textit{core\_course\_get\_contents}} que devuelve los módulos que hay en la página principal del curso y contiene un campo de tipo de módulo tal y como se puede ver en la figura \ref{fig:modname}.
	
	\imagenflotante{modname}{Extracto de la función del servicio web}
	\item Mediante el enlace al crear un módulo del curso. Figura \ref{fig:modnamelink}
	\imagenflotante{modnamelink}{Enlace que se genera mientras se crea el módulo}
	
\end{itemize}

\subsubsection{Tipos de módulos del curso}

En la tabla \ref{tabla:herramientasportipodeuso} muestra los diferentes tipos de módulos que se han encontrado en Moodle.


\tablaSmall{Herramientas y tecnologías utilizadas en cada parte del proyecto}{l c c c c}{herramientasportipodeuso}
{ \multicolumn{1}{l}{Código nemotécnico} & Nombre en español & Nombre en inglés \\}{ 
assign             & Tarea                    & Assignment          \\
book               & Libro                    & Book                \\
chat               & Chat                     & Chat                \\
choice             & Consulta                 & Choice              \\
data               & Base de datos            & Database            \\
lti                & Herramienta externa      & External tool       \\
feedback           & Encuesta                 & Feedback            \\
resource           & Archivo                  & File                \\
folder             & Carpeta                  & Folder              \\
forum              & Foro                     & Forum               \\
glossary           & Glosario                 & Glossary            \\
hotpot             & HotPot                   & HotPot              \\
imscp              & Paquete de contenido IMS & IMS content package \\
journal            & Diario                   & Journal             \\
label              & Etiqueta                 & Label               \\
lesson             & Lección                  & Lesson              \\
page               & Página                   & Page                \\
quiz               & Cuestionario             & Quiz                \\
scorm              & Paquete SCORM            & SCORM package       \\
survey             & Encuestas predefinidas   & Survey              \\
url                & URL                      & URL                 \\
wiki               & Wiki                     & Wiki                \\
workshop           & Taller                   & Workshop            \\
} 




\subsection{Ítem de calificación}
El Ítem de calificación\cite{noauthor_grade_nodate} (en inglés \textit{Grade item}), es cada fila de la tabla de calificaciones.
Las carpetas de la tabla, a parte de Ítem de calificación, se considera también como categoría (\textit{category}).

Los Ítem de calificación se puede mostrar mediante una tabla para un solo alumno (ver figura \ref{fig:item_calificacion}) o una tabla de todos los alumnos y todos los Ítems de calificaciones (ver figura \ref{fig:alternativa_grade_item})
	

\imagenflotante{item_calificacion}{Ejemplo de una tabla con varios Ítem de calificación.}

\imagenflotante{alternativa_grade_item}{Ejemplo de tabla parcial con los usuarios e Ítems de calificación}
\subsection{Categorías de cursos}
Las categorías de cursos\cite{noauthor_course_nodate-1} (en inglés \textit{course categories}) organizan los cursos de la página de Moodle de forma jerárquica (ejemplo en la figura \ref{fig:course_categories})
\imagenflotante{course_categories}{Ejemplo de categorías del curso enmarcados en rojo.}

\subsection{Registros}
Los registros es un reporte de actividades realizadas por los usuarios, existe registros de sitio web y del curso. En este proyecto solo nos centraremos en el segundo caso, estos se acceden en desde el curso en la rueda de configuración \textit{Aún más... ->  Informes ->  Registros -> Conseguir estos registros}

\imagen{registros}{Ejemplo de varias entradas de un registro.}

\subsubsection{Columnas de los registros}
Todas las entradas del registros se muestra en una tabla con varias columnas: 
\begin{enumerate}
	\item \textbf{Hora}: La fecha y hora.
	\item \textbf{Nombre completo del usuario}: Nombre y apellidos del usuario que realiza el evento o acción.
	\item \textbf{Usuario afectado}: Nombre y apellidos del usuario afectado por el evento o acción.
	\item \textbf{Contexto del evento}: En qué lugar del curso se ha producido el registro.
	\item \textbf{Componente}: puede indicar una actividad, un recurso o el sistema.
	\item \textbf{Nombre del evento}: que evento o acción se ha realizado.
	\item \textbf{Descripción}: contiene la información de las anteriores columnas pero mostrando su identificador único (\textit{id}) en vez de los nombres usuarios. Cada descripción es única en función del \textbf{Componente} y \textbf{Nombre de evento}.
	\item \textbf{Origen}: indica el tipo de fuente donde se ha generado el registro. Pueden ser:
	\begin{itemize}
		\item Servicio web (ws): servicios web de Moodle.
		\item Web: página web.
		\item CLI: desde el servidor.
	\end{itemize}
	\item \textbf{Dirección IP}: la dirección IP del dispositivo que ha activado el registro.
	
\end{enumerate}

\subsubsection{Opciones y descarga}
También ofrece la posibilidad de filtrar mediante diferente opciones y su descarga en varios formatos de texto (CSV, Excel, HTML, JSON, etc).

\subsubsection{Problema de visualización}

El problema de los registros que ofrece Moodle es la laxa forma de visualizar los datos a pesar de las opciones de filtrado. Sobre todo la información de las \textbf{Descripciones} es confusa al mostrar los id's de usuarios, módulos del curso, etc...





\capitulo{4}{Técnicas y herramientas}

En este apartado se explicarán las técnicas y herramientas usadas para el desarrollo del proyecto.

\section{Metodología SCRUM}

Se ha seguido una de las metodologías ágiles conocida como SCRUM adaptada para una sola persona con reuniones semanales en vez de diarias. Con una estrategia iterativa (incremental) con entregas parciales en los \textit{sprints}.

\section{Moodle}
\href{https://moodle.org}{Moodle} es una herramienta de gestión del aprendizaje y de distribución libre. Está desarrollado principalmente en Linux usando Apache, PostgreSQL/MySQL/MariaDB y PHP (también conocida como plataforma LAMP) \cite{noauthor_35/instalacion_nodate}.

Ayuda a los docentes en la educación a distancia y en la gestión de los diferentes modos de calificación.




En la página oficial de Moodle ofrece varias versiones de prueba que cualquiera puede acceder con unos usuarios y contraseña predefinidos.

\section{REST API}
REST \cite{noauthor_api_nodate} son todas las interfaces entre sistemas que usa HTTP para conseguir datos o realizar operaciones sobre esos datos en todos los formatos posibles, como puede ser XML y JSON.

En Moodle la funcionalidad de REST API se conoce como \textit{Web Services} que a partir de un \textit{token} generado previamente por la función \textbf{\textit{moodle\_mobile\_app}} permite el acceso a otras funciones del servicio. Puede ser desde ver las calificaciones de alumnos hasta crear actividades o tareas.

\section{Java}
Java es un lenguaje de programación orientado a objetos, diseñado para escribir el código una vez y ejecutar en cualquier otro dispositivo sin tener que recompilar siempre y cuando se tenga instalado\textit{Java Runtime Enviroment (JRE)}. Actualmente es uno de los lenguajes programación con mayor popularidad de uso, tanto en el ámbito de desarrollo como en el de usuarios.

El desarrollo, al igual que el anterior proyecto UBUGrades y UBULogs, se ha seguido con Java 8 al ser una edición con soporte extendido \cite{noauthor_oracle_nodate}.

\section{JavaFX}
Es un paquete gráfico que permite crear aplicación interactivas. Viene incluido en Java 8 dentro de la \textit{Java Development Kit (JDK)} y en \textit{Java Runtime Enviroment (JRE)} por lo que se puede ejecutar en otro dispositivo sin tener que instalar librerías adicionales.

\subsection{Scene Builder}
Es una herramienta que permite diseñar de una forma sencilla ventanas gráficas sin programar código. Los elementos de la ventana se añaden arrastrando y soltando (\textit{Drag and drop}) dentro de la herramienta.

\imagen{SceneBuilder}{Herramienta de Scene Builder}

Con Scene Builder generará un archivo FXML, muy parecido al formato XML, donde se guarda las posiciones de los elementos de la ventana gráfica. Desde Java solo tiene que cargar el fichero FXML.

\section{JavaScript}
Lenguaje de programación orientado a páginas web, usado principalmente en el lado del cliente. Al ser compatible con todos los navegadores web del mercado, se ha vuelto un estándar de facto.

En el proyecto se sigue usando Charts.js para la visualización las nuevas gráficas.

\section{CodePen}
\href{https://codepen.io}{CodePen} es una herramienta web que permite crear prototipos de páginas HTML, CSS y JavaScript y ver los resultados en tiempo real.

Usado para probar prototipos de gráficas con los registros del curso.

\section{JSON}

\textit{JavaScript Object Notation (JSON)}, formato de texto sencillo para el intercambio de datos. En principio se usaba como subconjunto de notación en JavaScript pero debido a su masiva popularidad en otros ámbitos, se considera hoy en día como lenguaje independiente.

Los datos de Moodle que se recibe mediante REST API está en formato de texto JSON.

\section{Eclipse}

\href{https://www.eclipse.org/}{Eclipse} es uno de los entornos de desarrollo (IDE) más famosos, junto con
\href{https://www.jetbrains.com}{\textit{IntelliJ IDEA}} y \href{https://netbeans.org/}{\textit{NetBeans}}, para Java. De código abierto con la capacidad de añadir nuevas características mediante plugins.

Se ha decidido emplear Eclipse frente a las otras alternativas por comodidad y es la herramienta que se maneja durante la carrera.

\section{GitHub}
Es una plataforma de desarrollo colaborativo para alojar proyectos usando el sistema de control de versiones Git. UBUGrades está alojado en GitHub.

Adicionalmente se ha usado la extensión del navegador ZenHub que se integra en GitHub y permite el uso de \textit{canvas} para la organización de los \textit{issues}.


\section{LaTeX}

Sistema de composición de textos dirigido a la creación de documentos con una buena calidad tipográfica. Actualmente muy popular en el ámbito de la investigación con artículos y libros científicos.

El presente proyecto está realizado en \LaTeX usando las versiones portables de MiKTeX y TeXstudio. 

\subsection{MiKTeX}

\href{https://miktex.org/}{MiKTeX} es una distribución de \LaTeX diferentes plataformas de sistemas operativos, es de código libre y viene incluido con múltiples paquetes de tipografías y macros.

\subsection{TeXstudio}

\href{https://www.texstudio.org/}{TeXstudio} es un entorno de desarrollo integrado (IDE) de edición de texto, código abierto con resaltado de sintaxis orientado a la creación de documentos \LaTeX.

\section{Zotero}
\href{https://www.zotero.org/}{Zotero} es un gestor de referencias bibliográficas de código libre creado por el Center for History and New Media de la Universidad George Mason.

\section{Python}
\href{https://www.python.org/}{Python} es un lenguaje de programación interpretado, orientado a objetos, dinámicamente tipado y de alto nivel. Este lenguaje es muy atractivo para crear prototipos de forma sencilla y rápida. 

Se ha empleado Python con el IDE \href{https://jupyter.org/}{Jupyter Notebook} programando varias funciones de utilidad para el programador. Una de estas funcionalidades genera código Java automáticamente.

\section{Dependencias externas de Java}
En este apartado se explicarán las librerías externas de Java usadas durante el desarrollo del proyecto. 

\subsection{Apache Commons CSV}
\href{https://commons.apache.org/proper/commons-csv/
}{Apache Commons CSV 1.6} permite la lectura y/o escritura de textos en \textit{comma-separated values} comúnmente conocido como formato \textbf{CSV}. 
El proyecto lo emplea para separar cada uno de las columnas y sus elementos de los registros del curso.

\subsection{Apache Commons Math}
\href{https://commons.apache.org/proper/commons-math/}{Apache Commons Math
3.6.1} es una librería orientado a las matemáticas y estadísticas.
Nos posibilita sacar datos estadísticos como las medias, desviaciones, máximos y mínimos solo enviando los valores.

Utilizado en el proyecto para las medias del curso tanto en las calificaciones como los registros.

\subsection{Gson}
\href{https://github.com/google/gson}{Gson 2.8.5} ofrece la importación y exportación de textos JSON a partir de los métodos \textit{toJson()} y \textit{fromJson()}.

Ignorado de momento para el proyecto. A implementar en el futuro para guardar los datos de los registros de todos los usuarios.

\subsection{SLF4J y Log4j}
Ambas librerías está dedicado a registrar en fichero y/o consola las acciones realizadas durante la ejecución del programa. De gran utilidad para encontrar errores del programa. 
Hay que tener en cuenta que \href{https://www.slf4j.org/}{SLF4J} es la capa de abstracción y \href{https://logging.apache.org/log4j}{Log4j} la implementación.

\subsection{jsoup}
\href{https://jsoup.org/}{jsoup 1.11.3} es una librería que trabaja con textos HTML, desde la descarga de páginas web hasta su tratamiento y extracción de datos de manera sencilla y cómoda.

Ha resultado ser realmente útil, muy polivalente en las funcionalidades. Empleado para la descarga de los registros, las imágenes de los usuarios y analizar el contenido en una de las funciones de Moodle que devuelve como HTML.

\subsection{ThreeTen Extra}
\href{https://www.threeten.org/threeten-extra/}{ThreeTen Extra 1.5.0} proporciona varias formas de medir el tiempo no disponibles en el paquete de Java Time. Las clases de ThreeTen Extra son complementarias y extienden de Java Time.

Ejemplos aplicados en el proyecto han sido la combinación de número de semana y año, o trimestre y año.

\section{Librerías de JavaScript}

Al manejar una ventana de navegador web para mostrar las gráficas, se han apoyado por varias librerías de JavaScript.

\subsection{Chart.js}
\href{https://www.chartjs.org/}{Chart.js} es una librería de de visualización de datos convertidos en múltiples tipos de gráficas.

Se ha mantenido está librería por su facilidad de crear las gráficas y sus múltiples opciones. Sin embargo por un error de la librería en el anterior proyecto, se ha actualizado Chart.js a la última versión disponible \href{https://www.chartjs.org/dist/2.8.0/Chart.min.js}{v2.8.0}.

\subsection{Google Charts}

\href{https://developers.google.com/chart/}{Google Charts} es una librería gratuita de gráficas muy fáciles de usar, dispone de múltiples tipos de gráficas interactivas. En sus términos y condiciones de uso no permite descargar el fichero JavaScript y usarlo en local \cite{noauthor_frequently_nodate}. En vez de eso, hay que utilizar el enlace que proporciona Google \cite{noauthor_quick_nodate}.

Se usa para mostrar la tabla de calificaciones de los alumnos. 

\subsection{color-hash}

color-hash \cite{zeng_generate_2019} es una librería que permite generar colores en función de la cadena de texto usado.

\subsection{html2canvas}

\href{https://html2canvas.hertzen.com/}{html2canvas}
Realiza capturas de pantalla en una página web con JavaScript.



\section{Sonarqube y Codacy}

\href{https://www.sonarqube.org/}{SonarQube} y \href{https://www.codacy.com/}{Codacy} son herramientas de evaluación del código fuente, muestra los códigos duplicado, errores, vulnerabilidades y malas prácticas de programación. 
Se instaló el plugin de Eclipse de Sonarqube llamado \href{https://www.sonarlint.org/}{SonarLint}.

\capitulo{5}{Aspectos relevantes del desarrollo del proyecto}

En este apartado se comenta los aspectos importantes que han surgido durante el desarrollo del proyecto. Las decisiones tomadas relativas a la funcionalidad de la aplicación y qué consecuencias ha tenido. También explicaremos sobre los problemas encontrados y las soluciones aplicadas.

\section{Sin acceso a permisos de profesor en el Moodle de la Universidad (UBUVirtual)}

Uno de los principales problemas para realizar las pruebas no se ha podido tener acceso a asignaturas con rol de profesor. Es decir, muchas funciones de Moodle no son accesibles con el rol de estudiante mientras que los profesores si lo tienen.


Se ha instalado una versión en el equipo para realizar las pruebas, se comentara con más detalle en la sección \ref{sec:instalacionMoodle}. También se ha empleado una versión de prueba oficial del propio equipo de desarrollo \href{https://school.demo.moodle.net}{\textit{Mount Orange School}} en la versión \textbf{3.7}, probablemente se use esta versión en el aula virtual de la Universidad de Burgos en el futuro.

Las comprobaciones en un entorno real lo han realizado varios profesores docentes de la universidad, entre ellos los tutores, en UBUVirtual. 

Cabe destacar que los tutores probaron cada mejora implementada en un corto lapso de tiempo, es decir en menos de 24 horas laborales ya tenía el visto bueno o errores encontrados. Esto ha facilitado mucho los tiempos de desarrollo de nuevas funcionalidades.

\section{Instalación de la versión limpia de Moodle} \label{sec:instalacionMoodle}
Después de conocer la versión de Moodle que usa la Universidad (\textbf{3.5.1}) mediante  la función \textbf{\textit{core\_webservice\_get\_site\_info}}, se ha intentado instalar en mi equipo la misma versión pero ya no estaba disponible en la página de descargas \cite{noauthor_moodle_nodate}. Finalmente se ha usado la \textbf{3.5.4+}.

A pesar de usar una versión con cambios menores respecto al de la Universidad. Si que se encontró una diferencia reseñable. 

\subsection{Token para el ingreso} \label{sec:logintoken}
Añaden el Token para ingreso\cite{noauthor_token_nodate} (en inglés \textit{login token}), una característica relacionada con la seguridad introducida en las versiones de Moodle 3.1.15, 3.3.9, 3.4.6, 3.5.3 y 3.6.0. Ayuda a proteger frente a vulnerabilidades como el robo de sesión de los usuarios. Este token de ingreso aparece como un \textit{input} HTML oculto\cite{noauthor_html_nodate-1} al iniciar sesión y se envía el formulario de ingreso junto con el usuario y contraseña.

\imagen{logintoken}{Ejemplo de Login Token en el inicio de sesión de la página \href{https://school.demo.moodle.net/login/index.php}{Mount Orange School}}.

Para solucionar este problema y que funcione en ambas versiones y futuras, cada vez que el usuario inicia sesión, se busca en la página el \textit{input logintoken} tal y como aparece en la figura \ref{fig:logintoken} con métodos de Web Scraping\cite{marti_que_2016}. Si lo encuentra manda el valor como parámetro adicional de la URL, en caso contrario no se manda este parámetro.

Hay que destacar que el Token de ingreso solo afecta si se quiere iniciar sesión en la página del servicio de Moodle (necesario para poder descargar los registros del curso), en ningún caso es necesario para los servicios Web.

\section{Empezar de cero la parte de los registros}

Se decidió prescindir totalmente del trabajo de UBULogs y empezar de nuevo debido a las modificaciones de Moodle como el token de ingreso (ver sección \ref{sec:logintoken}). 

Finalmente se ha seguido por una vía diferente al mencionado proyecto, descargar los CSV directamente en memoria sin guardarlo en un fichero temporal. También se ha aplicado una mejora para bajarse los registros de forma diaria cuando se actualiza los datos.

Otra mejora reseñable, al estar las horas de los registros mostradas en función de la zona horaria usada por el usuario, se ha decidido tomar como referencia para la aplicación la zona horaria del servidor. No importa si el usuario esta usando otra que no sea la misma que el servidor.

\section{Decisión sobre seguir el código fuente de UBUGrades o refactorizar}

Cuando pasaron dos semanas desde el inicio del proyecto, acabé en la tesitura sobre seguir con el código fuente o refactorizar la parte lógica de la aplicación (back-end) aplicando patrones de diseño y tenga mejor mantenibiliad. 

Dos de los temores de tomar la decisión de refactorizar fue quedarme a medio camino en la entrega del proyecto o problemas de integración con el vista e interacción con el usuario (front-end).

Finalmente, después de otras dos semanas de reflexión, se decide refactorizar el código de la parte del back-end completamente. Al empezar de nuevo se debe dedicar un esfuerzo extra en el proyecto.



\begin{itemize}
	\item version de ubuvirutal no disponible, usado como pruebas una version parecida
	\item debate sobre si trabajar con el codigo o empezar de 0 el back end, 2 semanas de reflexion. Apuesta arriesgada debido a quedarme a medio camino en la entrega o problemas de integración con el front end.
	\item usado mas tiempo en pensar como estructurar las clases de forma que sea facil de leer y facil de modificar, que programando
	\item retraso de tiempos al haber un error de moodle en una funcion que saltaba error si el feedback del calificador esta oculto, se ha tenido que usar otra funcion mas engorrosa de trabajar (HTML)
	\item dedicado varias semanas a elegir como mostrar la grafica de registros, probado con varios prototipos de graficas de barras diferentes, con y sin plugins para chartjs
	\item ademas de los prototipos de las barras tambien se ha estado pensando y probando como estructurar el codigo para que sea lo mas sencillo posible añadir nuevos tipos de agrupaciones de fechas, crear los datasets de las barras, que sea sencillo añadir para componentes, eventos
\end{itemize}


\capitulo{6}{Trabajos relacionados}

En cuanto a los trabajos relacionados tenemos dos precedentes comentados anteriormente, UBUGrades y UBULogs hechos como Trabajos de Fin de Grado por alumnos de la universidad.

Existe varios plugins de Moodle que realizan trabajos parecidos al proyecto.

\section{Heatmap}

Heatmap \cite{noauthor_moodle_nodate-2} superpone un mapa de calor (gradiente de color) en un curso para resaltar los módulos dependiendo del nivel de actividad y así ayudar a los docentes a mejorar sus cursos. Ejemplo en la figura \ref{fig:heatmap}.

\imagenflotante{heatmap}{Ejemplo de Heatmap}

\section{SmartKlass™ Learning Analytics Moodle}
SmartKlass™ Learning Analytics Moodle \cite{noauthor_moodle_nodate-3} es un panel de análisis dirigido a docentes y estudiantes que muestra la evolución del rendimiento sobre el curso. Ejemplo en la figura \ref{fig:smartklass}.

Aparentemente el proyecto lleva sin actualizarse desde hace 3 años\cite{noauthor_klassdata/moodle-local_smart_klass:_nodate}.

\imagenflotante{smartklass}{Ejemplo del panel SmartKlass}

\section{Engagement analytics}

Engagement analytics \cite{noauthor_moodle_nodate-4} proporciona información sobre el progreso de los estudiantes en relación con una serie de indicadores.

Como su nombre indica, analiza el nivel de compromiso o participación del estudiante en el curso.

\section{Course dedication}

Course dedication \cite{noauthor_moodle_nodate-5} Calcula el tiempo estimado de dedicación de los participantes en el curso.




\capitulo{7}{Conclusiones y Líneas de trabajo futuras}

Todo proyecto debe incluir las conclusiones que se derivan de su desarrollo. Éstas pueden ser de diferente índole, dependiendo de la tipología del proyecto, pero normalmente van a estar presentes un conjunto de conclusiones relacionadas con los resultados del proyecto y un conjunto de conclusiones técnicas. 
Además, resulta muy útil realizar un informe crítico indicando cómo se puede mejorar el proyecto, o cómo se puede continuar trabajando en la línea del proyecto realizado. 


\printbibliography

\end{document}
